% Chapter 1: Introduction
\chapter{Introduction}\label{ch:introduction}

%----------------------------------------------------------------------------
\section{Background}\label{sec:background}

\guide{Give the reader the facts they need to understand the topic and how it
links to your research question. Start from the broad theme, then narrow down
to the specific details that make your problem relevant. By the end, the need
for the research should already be clear.}

When people watch a video, they do not look everywhere. Their eyes move several
times per second, and different observers tend to look at the same few regions:
moving objects, faces and, in conversations, the person who is
speaking~\cite{mital2011clustering,coutrot2014how}. Where people look is shaped
by the image itself, by the viewer's task and by strong regularities such as
the preference for the centre of the
screen~\cite{itti2001computational,borji2013state,tseng2009quantifying}.
\emph{Video saliency prediction} is the task of predicting, for every frame of a
video, the distribution of human fixations. Models of this kind support
applications in which the allocation of attention matters, from compression and
summarisation to the design and evaluation of visual
content~\cite{borji2021saliency}, and they are also scientific models of human
attention whose successes and failures say something about what drives gaze.

The field has progressed in step with the representations it builds on. Early
models computed saliency from the local contrast of hand-crafted features such
as intensity, colour, orientation and motion~\cite{itti1998model}. The
introduction of large eye-tracking benchmarks, most importantly
DHF1K~\cite{wang2018revisiting,wang2021revisiting}, enabled deep networks to be
trained for the task, and the accuracy of these networks has risen with each
generation of their encoders: image networks with recurrent layers,
three-dimensional convolutional networks pretrained for action recognition,
video transformers and, most recently, video foundation models trained with
self-supervision on large collections of video~\cite{moradi2024salfom}.

Self-supervised video models differ in what they learn to predict. One family,
the joint-embedding predictive architectures, learns by predicting the
\emph{representations} of hidden parts of a video from the visible
parts~\cite{assran2023self,bardes2024revisiting}. Such models have two
properties that make them interesting for saliency prediction. Their features
are reported to work well without adapting the encoder to a new
task~\cite{bardes2024revisiting}, which suggests that a frozen encoder with a
small task-specific decoder might suffice. And they contain a trained predictor,
whose errors indicate what the model failed to anticipate, a learned analogue of
the \enquote{surprise} that has been proposed as a driver of human
attention~\cite{itti2009bayesian}.

%----------------------------------------------------------------------------
\section{Motivation}\label{sec:motivation}

\guide{Explain what motivated you to carry out this work.}

Three considerations motivated this work. The first is practical.
SalFoM~\cite{moradi2024salfom}, the first video saliency model built on a video
foundation encoder, which reported the best DHF1K test results among the methods
it compared with, fine-tunes a 300-million-parameter encoder on a few hundred
training videos and runs it once for every output frame. Fine-tuning
so large a model on so little data risks overfitting, which its authors offer as
an explanation for its weaker results on a small dataset, and the per-frame
design makes inference expensive. A frozen encoder with a decoder that predicts
many frames at once would avoid both problems, if its features are good enough.

The second consideration is scientific. The idea that attention is drawn to the
unexpected has strong support from computational models built on hand-crafted
features~\cite{itti2009bayesian}, but it had not been tested with the learned
predictor of a modern video model. If such a predictor's errors marked where
people look, they would provide a principled attention cue learned without any
gaze data; if they did not, understanding why would clarify what these models
actually learn.

The third consideration is methodological. Published comparisons in video
saliency rely mostly on numbers copied from other papers, computed with
different evaluation code and sometimes on different data splits, and they
rarely report uncertainty. Small reported gains are therefore hard to interpret.
A study that re-scores its main baseline with the same validated evaluator,
quantifies uncertainty and fixes its decision rules in advance can give answers
that are more reliable, including negative answers.

%----------------------------------------------------------------------------
\section{Statement of the Problem}\label{sec:problem}

\guide{Define the problem in the literature, theory or practice that this
thesis solves, clearly enough that the reader knows exactly what is being
addressed. A good problem statement answers \enquote{Why does this research
need to be conducted?}}

Recent video saliency models obtain their accuracy from large pretrained
encoders, which they adapt to small saliency datasets, at a high computational
cost per output frame. It is not known whether this adaptation is necessary, or
whether the representations of a self-supervised video encoder trained by
feature prediction are sufficient on their own, with only a light decoder
trained on gaze data. Nor is it known whether the prediction errors of such an
encoder's predictor, which offer a learned measure of surprise, carry
information about where people look. Finally, because comparisons in the field
are rarely made under a common evaluation protocol with uncertainty estimates,
it is unclear which design choices, beyond the encoder, still matter once a
strong encoder is used. This thesis addresses these questions with a frozen
V-JEPA~2 encoder~\cite{murlabadia2026vjepa21}, a purpose-built decoder, and a
pre-registered evaluation on nine eye-tracking test sets.

%----------------------------------------------------------------------------
\section{Research Questions}\label{sec:questions}

\guide{List the questions the thesis answers. \Cref{ch:conclusion} should
answer each one explicitly.}

This study answers the following research questions:
\begin{enumerate}[label=\textbf{RQ\arabic*.},leftmargin=*]
  \item Can a frozen self-supervised video encoder of the V-JEPA family, combined
    with a lightweight spatio-temporal decoder, predict human fixations in video
    more accurately than the closest foundation-model approach, SalFoM, when both
    are evaluated under an identical protocol?
  \item Does the prediction error of the encoder's pretrained latent predictor
    (\enquote{predictive surprise}) carry information about where people look,
    beyond that carried by simple appearance and motion cues?
  \item Which design factors determine accuracy when the encoder is frozen: the
    pre-training and scale of the encoder, parameter-efficient adaptation, the
    learned centre prior, deep supervision, static-image co-training, temporal
    sampling, the amount of training data and ensembling?
  \item How well does the approach transfer to eye-tracking datasets with other
    content and viewing tasks, and does adding the soundtrack improve it on
    audio-visual datasets?
\end{enumerate}

%----------------------------------------------------------------------------
\section{Objectives}\label{sec:objectives}

\guide{State a general objective, and specific objectives that are simple,
specific and follow from the problem statement.}

\subsection{General Objective}

The general objective of this study is to design, implement and rigorously
evaluate a video saliency prediction model built on a frozen self-supervised
video encoder, and to test whether the encoder's own prediction error is a
useful cue for human attention.

\subsection{Specific Objectives}

To achieve the general objective, the study will:
\begin{itemize}
  \item review the literature on visual attention, video saliency prediction and
    self-supervised video representation learning;
  \item prepare the eye-tracking datasets and an evaluation protocol whose metric
    implementation is validated against reference code, including a new
    ground truth for Hollywood-2 from openly licensed gaze recordings;
  \item design and implement a decoder that turns the intermediate features of a
    frozen V-JEPA~2 encoder into saliency maps for many frames per forward pass;
  \item test the predictive surprise of the V-JEPA~2 predictor as a fixation
    cue against simple reference cues;
  \item compare encoders and design choices under decision rules fixed before
    each experiment;
  \item evaluate the final models on DHF1K and eight further test sets, with
    paired statistical comparisons where the data allow; and
  \item measure the computational cost of the approach.
\end{itemize}

%----------------------------------------------------------------------------
\section{Research Methodology}\label{sec:methodology-overview}

\guide{A short overview of how the objectives were achieved: literature
review, data, tools and evaluation. Keep it brief and point to
\cref{ch:methodology} for the details.}

\subsection{Literature Review}
The literature was searched in publisher indexes and through the reference
lists of key papers, covering visual attention, image and video saliency
prediction, video representation learning and evaluation methodology. Only
peer-reviewed publications are cited. The venue and bibliographic details of
every reference were verified against the publishers' records or, for
proceedings without digital object identifiers, against the indexes of those
proceedings; works available only as preprints were excluded.

\subsection{Data Collection}
No new gaze data were recorded. The study uses public eye-tracking datasets:
DHF1K, UCF Sports, Hollywood-2 with the gaze recordings of Vig et al., the six
audio-visual datasets of the STAViS benchmark, and SALICON for one ablation
(\cref{sec:dataset}).

\subsection{Tools and Techniques}
The model was implemented in Python with PyTorch~\cite{paszke2019pytorch}; the
pretrained encoders were loaded with the Transformers
library~\cite{wolf2020transformers}. All experiments ran on a single
workstation with one NVIDIA RTX A5000 graphics card (\cref{sec:setup}).

\subsection{Evaluation}
Models were evaluated with six standard saliency metrics~\cite{bylinskii2019metrics},
computed by an evaluator validated against reference code. Differences between
models were tested with a paired bootstrap over videos, and every decision was
governed by a rule written before the corresponding experiment
(\cref{sec:design}).

%----------------------------------------------------------------------------
\section{Scope and Limitations}\label{sec:scope}

\guide{State what the study covers, and what it does not cover and why.}

\paragraph{Scope.} The study covers offline (non-causal) saliency prediction for
video, in which the map of a frame may depend on the frames around it. The main
model is visual; the soundtrack is examined in one separate experiment. The
encoders are frozen V-JEPA~2 models in two sizes; V-JEPA~2.1 and UMT-L are
included in the encoder comparison. Evaluation covers nine test sets: DHF1K,
UCF Sports, Hollywood-2 with the Vig gaze data, and six audio-visual datasets.

\paragraph{Limitations.} The DHF1K test-set annotations are withheld by the
dataset authors; predictions for the test set were prepared, but their scores
were not available when the thesis was written, so DHF1K results are reported on
the official validation set. The standard Hollywood-2 ground truth could not be
obtained, so the Hollywood-2 results use other gaze data and are not comparable
with published numbers. Full fine-tuning of the encoder, other foundation
encoders and causal (streaming) prediction were outside the scope of the study,
and the component ablations used a single training seed because of the
available computing resources.

%----------------------------------------------------------------------------
\section{Contributions}\label{sec:contributions}

\guide{List the concrete contributions of the thesis (new method, dataset,
results, analysis) and how they extend existing knowledge.}

The main contributions of this thesis are:
\begin{itemize}
  \item \textbf{SURGE-D}, a video saliency model that combines a frozen V-JEPA~2
    encoder with a light feature-pyramid decoder of factorised spatio-temporal
    convolutions and predicts maps for sixteen frames per forward pass. With
    7.4 million trained parameters, its three-seed ViT-g ensemble reaches a
    correlation coefficient of 0.603 on the DHF1K validation set, against 0.564
    for SalFoM scored by the same evaluator, and is better on all six metrics
    with paired bootstrap intervals that exclude zero. Counting the encoder
    alone, SalFoM's per-frame design requires about six times more computation
    per output frame.
  \item \textbf{A negative result on predictive surprise}: the error of the
    pretrained V-JEPA~2 predictor does not indicate where people look; it is
    beaten by frame difference and is even anti-correlated with fixations,
    because it mainly reflects the atypicality of regions rather than violated
    expectations.
  \item \textbf{Controlled evidence on what matters} once the encoder is frozen:
    the encoder's pre-training and scale produce the only large gains, seed
    ensembles a small one, while low-rank adaptation, the learned centre prior,
    deep supervision, static-image co-training, longer clips, more training
    videos and a simple audio fusion produce none that could be measured.
  \item \textbf{A transparent evaluation} across nine test sets with a validated
    evaluator, pre-registered decision rules, paired statistics and a new,
    openly reproducible Hollywood-2 ground truth from the gaze recordings of Vig
    et al., with the code, splits and decision log provided with the thesis.
\end{itemize}

%----------------------------------------------------------------------------
\section{Organization of the Thesis}\label{sec:organization}

The rest of this thesis is organized as follows.
\Cref{ch:literature} reviews the background concepts and techniques.
\Cref{ch:related} discusses related work and identifies the research gap.
\Cref{ch:methodology} presents the proposed approach and the evaluation protocol.
\Cref{ch:experiments} describes the experiments and discusses the results.
\Cref{ch:conclusion} concludes the thesis and outlines future work.
